Let $\alpha=\angle BAC=\angle CBE$. Since $AD=DE$, triangle ADE has angle AED equal to alpha, so $\angle BDE=2\alpha$. Since $DE=BE$, $\angle DBE=2\alpha$, giving $\angle ABC=3\alpha$. Since $BE=BC$, $\angle BCE=\dfrac{(180^\circ-\alpha)}{2}$. Summing ABC angles gives $\alpha+3\alpha+\dfrac{(180^\circ-\alpha)}{2}=180^\circ$, hence $\alpha=\dfrac{180^\circ}{7}$. Therefore $\angle BED=180^\circ-4\alpha=\boxed{\dfrac{540^\circ}{7}}$.
\begin{asy}
import graph;import geometry;usepackage("amsmath");size(7cm);
real a=180.0/7;pair A=(0,0),B=(2*cos(2*pi/7)+1,0),D=(1,0),E=intersectionpoint(A--4*dir(a),D--D+4*dir(180-2*a)),C=intersectionpoint(A--8*dir(a),B--B+8*dir(180-3*a));draw(A--B--C--cycle);draw(D--E^^B--E,blue);label("$A$",A,(-1,-1));label("$B$",B,(1,-1));label("$C$",C,(0,1));label("$D$",D,(0,-1));label("$E$",E,(-1,1));
\end{asy}

